\section{Preliminaries}
\label{sec:background}
\subsection{Write Ahead Logging}


\subsection{P-WAL}


% Problem of P-WAL
P-WAL is attractive because it removes the single-WAL insertion bottleneck by assigning transactions to multiple log streams.  However, combining P-WAL with SSN exposes three
problems.

\textbf{P1. Multiple Global Counter Accesses:~} First, P-WAL can introduce frequent global-counter accesses.  Even if the CC
layer already assigns a commit timestamp for SSN validation and serialization, the recovery layer may allocate a separate global LSN or maintain a global durable prefix.  This duplicates the ordering work already performed by the CC layer and adds shared-memory synchronization to the commit path.

\textbf{P2. Waiting for Data Transfer:~} Second, workers can spend a long time waiting around \texttt{fdatasync}.  P-WAL parallelizes log streams, but a transaction cannot be acknowledged until the corresponding durable condition is satisfied.  If worker threads wait directly
for log flushes or durable-prefix updates, storage synchronization and notification latency reduce the effective scalability of the CC protocol.

\textbf{P3. Total History:~} Third, global-prefix acknowledgment waits for log records that are unrelated to the transaction being acknowledged.  This rule is safe, but overly conservative:
if one WAL shard is slow, transactions on other shards may be delayed even when they have no dependency on transactions logged by the slow shard.  The opposite design point, local-only acknowledgment, avoids this delay but is unsafe.  If transaction $T$ reads a value written by transaction $U$, acknowledging $T$ before $U$ is recoverable may allow a crash recovery state that contains $T$ without the state on which $T$ depends.
